%==================================================================

\begin{defn}[пучок $\Shv{F}^{\#}$]
\begin{equation}\label{eq:Fsharp}
 \Shv{F}^{\#}(U):=\bigl\{\,s\colon U\to\et(\Shv{F})\ \bigm|\ s
 \text{ непрерывно},\ \pi\circ s=\id_U\,\bigr\}.
\end{equation}
Отображения ограничения --- обычные ограничения функций: при
$V\subset U$ отображение $\Shv{F}^{\#}(U)\to\Shv{F}^{\#}(V)$ ---
$s\mapsto s|_{V}$.
\end{defn}

Разберём два условия в~\eqref{eq:Fsharp}:
\begin{conditions}
\item $\pi\circ s=\id_U$ означает, что $s(x)$ лежит над $x$, т.\,е.
 $s(x)\in\Shv{F}_x$: сечение выбирает росток в каждой точке.
\item Непрерывность означает (по теореме~\ref{thm:open}), что для любого
 открытого $G\subset\et(\Shv{F})$ множество $\{x\in U: s(x)\in G\}$
 открыто. Это и есть <<плавность>> выбора.
\end{conditions}

В обеих записях встречается уточнение <<$\pi\circ s=\id_U$», а не
<<$(\pi\circ s)(U)=\id_U$»: имеется в виду тождественность как
отображений, т.\,е. равенство в каждой точке.

\begin{rem}[сравнение с $\Shv{F}(U)$]
Для предпучка определено каноническое отображение
\begin{equation}\label{eq:i}
 i\colon\Shv{F}(U)\longrightarrow\Shv{F}^{\#}(U),\qquad
 i(s)=\tilde s,\quad \tilde s(x)=s_x.
\end{equation}
Предпучок <<теряет>> часть сечений (отождествляет те, что локально
совпадают), а пучок <<добавляет>> (склеивает согласованные локальные
наборы). Между этим и стоит весь разговор о прокачивании.
\end{rem}

\begin{bookbox}
\booksrc{Годеман, Замечание 1.2.3, с.~133: каждый предпучок
 $\Shv{F}$ канонически определяет пучок $\tilde{\Shv{F}}$,
 \emph{порождённый предпучком $\Shv{F}$}, и канонический гомоморфизм
 предпучков $\Shv{F}\to\tilde{\Shv{F}}$ --- это и есть наше
 $i\colon\Shv{F}\to\Shv{F}^{\#}$. Годеман, гл.~II, \S\,1, п.~1.4,
 с.~134: <<простые пучки>> --- частный случай описанной конструкции.
}
\end{bookbox}

%==================================================================
